\section{Data and model}
\label{sec:model}

\paragraph{Catalogue.} We select every Gaia DR3 source with parallax
$\varpi>20$~mas (inside \CatRadius~pc) and $\varpi/\sigma_\varpi\ge10$, and keep
only members of the Gaia Catalogue of Nearby Stars (GCNS)
\cite{smart2021gcns}. This step matters. About 18\% of the raw parallax cut is
not in the GCNS, and those rows have a median tangential speed of 2.6~km/s,
which is the signature of distant background sources with spurious parallaxes.
With them included, star density rises with distance and stars appear to fall
toward the Sun. After the cut the density is flat at \Density~pc$^{-3}$ in
every 10-pc shell. Four bright stars that saturate Gaia ($\alpha$~Cen~A/B,
Sirius~A, Procyon~A) and the two nearest brown dwarfs (Luhman~16AB
\cite{luhman2013discovery} and WISE~0855$-$0714 \cite{luhman2014discovery}) are
added from the literature. The result is \NStars\ objects. \PctRV\% have a
radial velocity, from Gaia DR3 or the GCNS literature compilation.
\NHosts\ are confirmed exoplanet hosts (NASA Exoplanet Archive, matched by
Gaia DR3 identifier).

\paragraph{Missing radial velocities.} We fit a Gaussian velocity ellipsoid
$\mathcal N(\mu,\Sigma)$ to stars with RVs and draw each missing RV from that
Gaussian \emph{conditioned on the star's measured tangential velocity}. Drawing
RVs unconditionally breaks the coherence of the velocity field and makes faint
stars converge on the Sun. In our tests that error tripled the star count
inside 15~pc within 1~Myr. A unit test checks that the conditional fill is
unbiased on a synthetic field.

\paragraph{Kinematics.} Positions (pc) and velocities (pc/Myr) follow from
standard astrometric formulae and agree with \texttt{astropy} to $10^{-10}$.
Stars move on straight lines. Over $\lesssim1$~Myr inside 50~pc, Galactic
curvature bends paths by far less than the inter-star spacing, the same regime
used for Gaia encounter studies. With linear motion, Gliese~710 passes
the Sun at 0.051~pc (10,600~au) in 1.29~Myr, in line with the Gaia DR3
encounter analysis \cite{bailerjones2022gl710}.

\paragraph{Ships.} A ship refuels at every star it visits, unless stated
otherwise, and has a maximum single-hop range $r$. Range is the natural single
number because for a coasting ship it is the product of cruise speed and the
time the craft must survive in transit (Section~\ref{sec:engineering}). We
score stars within an evaluation radius well inside the catalogue: \EvalRadius~pc
for frozen stars and \TdEval~pc for moving stars. The outer shell acts as a
boundary buffer.
